\documentclass[12pt]{article}

\usepackage{amssymb,amsthm,enumitem, mathtools}
\usepackage{placeins}
\usepackage{graphicx}
\usepackage{subcaption}
\usepackage[mathscr]{euscript}
\graphicspath{ {./images/} }
\usepackage[margin=0.75in]{geometry}
\usepackage{multicol}
\usepackage{fancyhdr}
\usepackage[export]{adjustbox}
\usepackage{siunitx}
\usepackage{xcolor}
\usepackage{listings}
\sisetup{
  exponent-mode = scientific, % force scientific notation when needed
  exponent-base = 10,
  exponent-product = \times
  % IMPORTANT: do not set output-exponent-marker here
}

\pagestyle{fancy}
\fancyhf{}
\lhead{Horacio Moreno Monta\~nes}
\rhead{ME 523}
\rfoot{Page \thepage}
\setlength{\headheight}{15pt}

% theorems
\usepackage{thmtools}
\usepackage[framemethod=TikZ]{mdframed}
\mdfsetup{skipabove=1em,skipbelow=0em, innertopmargin=5pt, innerbottommargin=6pt}

\theoremstyle{definition}

\makeatletter

\declaretheoremstyle[headfont=\bfseries\sffamily, bodyfont=\normalfont, mdframed={ nobreak } ]{thmgreenbox}
\declaretheoremstyle[headfont=\bfseries\sffamily, bodyfont=\normalfont, mdframed={ nobreak } ]{thmredbox}
\declaretheoremstyle[headfont=\bfseries\sffamily, bodyfont=\normalfont]{thmbluebox}
\declaretheoremstyle[headfont=\bfseries\sffamily, bodyfont=\normalfont]{thmblueline}
\declaretheoremstyle[headfont=\bfseries\sffamily, bodyfont=\normalfont, numbered=no, mdframed={ rightline=false, topline=false, bottomline=false, }, qed=\qedsymbol ]{thmproofbox}
\declaretheoremstyle[headfont=\bfseries\sffamily, bodyfont=\normalfont, numbered=no, mdframed={ nobreak, rightline=false, topline=false, bottomline=false } ]{thmexplanationbox}


\declaretheorem[numberwithin=section, style=thmgreenbox, name=Definition]{definition}
\declaretheorem[sibling=definition, style=thmredbox, name=Corollary]{corollary}
\declaretheorem[sibling=definition, style=thmredbox, name=Proposition]{prop}
\declaretheorem[sibling=definition, style=thmredbox, name=Theorem, numbered=no]{theorem}
\declaretheorem[sibling=definition, style=thmredbox, name=Lemma, numbered=no]{lemma}



\declaretheorem[numbered=no, style=thmexplanationbox, name=Proof]{explanation}
\declaretheorem[numbered=no, style=thmproofbox, name=Proof]{replacementproof}
\declaretheorem[style=thmbluebox,  numbered=no, name=Exercise]{ex}
\declaretheorem[style=thmbluebox,  numbered=no, name=Example]{eg}
\declaretheorem[style=thmblueline, numbered=no, name=Remark]{remark}
\declaretheorem[style=thmblueline, numbered=no, name=Note]{note}
\declaretheorem[style=thmblueline, numbered=no, name=Recall]{recall}

\renewenvironment{proof}[1][\proofname]{\begin{replacementproof}}{\end{replacementproof}}



\AtEndEnvironment{eg}{\null\hfill$\diamond$}%

\newtheorem*{uovt}{UOVT}
\newtheorem*{notation}{Notation}
\newtheorem*{previouslyseen}{As previously seen}
\newtheorem*{problem}{Problem}
\newtheorem*{observe}{Observe}
\newtheorem*{property}{Property}
\newtheorem*{intuition}{Intuition}


% You can make short-forms for commands you want to use often.
% Here is an example that makes the integer Z symbol.
% In the text you just write '\ZZ' and it will make the symbol.\
\newcommand{\NN}{\mathbb{N}}
\newcommand{\ZZ}{\mathbb{Z}}
\newcommand{\RR}{\mathbb{R}}
\newcommand{\QQ}{\mathbb{Q}}
\newcommand{\KK}{\mathbb{K}}
\newcommand{\TT}{\mathscr{T}}
\newcommand{\PP}{\mathbb{P}}
\newcommand{\UU}{\mathscr{U}}
\newcommand{\OO}{\mathscr{O}}
\newcommand{\CC}{\mathbb{C}}
\newcommand{\DD}{\mathcal{D}}
\newcommand{\FF}{\mathcal{F}}
\newcommand{\BB}{\mathcal{B}}
\newcommand{\EE}{\mathbb{E}}
\newcommand{\mAA}{\mathcal{A}}
\newcommand{\Si}{\mathrm{Si}}
\newcommand{\Ai}{\mathrm{Ai}}
\newcommand{\Bi}{\mathrm{Bi}}
\newcommand{\Erf}{\mathrm{Erf}}
\newcommand{\erf}{\mathrm{erf}}
\newcommand{\im}{\textrm{im}}
\newcommand{\inv}[1]{#1^{-1}}
\newcommand{\cyclic}[1]{\ < \negmedspace #1 \negmedspace > \ }  
\newcommand{\rad}[1]{\textrm{rad($#1$)}}
\newcommand{\st}{\text{ s.t. }}
\newcommand*{\QED}{\null\nobreak\hfill\ensuremath{\square}}%
\renewcommand{\leq }{\leqslant}
\renewcommand{\geq}{\geqslant}
%\renewcommand{\phi}{\varphi}
\renewcommand{\epsilon}{\varepsilon}

\newcommand*\tageq{\refstepcounter{equation}\tag{\theequation}}

\newcommand{\dd}{\mathrm{d}}
\newcommand{\ee}{\mathrm{e}}
\newcommand{\ii}{\mathrm{i}}

\newcommand*\oline[1]{%
  \vbox{%
    \hrule height 0.5pt%                  % Line above with certain width
    \kern0.25ex%                          % Distance between line and content
    \hbox{%
      \kern-0.1em%                        % Distance between content and left side of box, negative values for lines shorter than content
      \ifmmode#1\else\ensuremath{#1}\fi%  % The content, typeset in dependence of mode
      \kern-0.1em%                        % Distance between content and left side of box, negative values for lines shorter than content
    }% end of hbox
  }% end of vbox
}

\newenvironment{sysmatrix}[1]
 {\left[\begin{array}{@{}#1@{}}}
 {\end{array}\right]}
\newcommand{\ro}[1]{%
  \xrightarrow{\mathmakebox[\rowidth]{#1}}%
}
\newlength{\rowidth}% row operation width
\AtBeginDocument{\setlength{\rowidth}{3em}}
%\setlength{\parindent}{0pt} 

\newcommand{\Res}[2]{\text{Res} \left ( #1 ; z = #2 \right)}
\newcommand{\Resw}[2]{\text{Res} \left ( #1 ; w = #2 \right)}

\usepackage{xifthen}

\def\testdateparts#1{\dateparts#1\relax}
\def\dateparts#1 #2 #3 #4 #5\relax{
    \marginpar{\small\textsf{\mbox{#1 #2 #3 #5}}}
}

\def\@lesson{}%
\newcommand{\lesson}[2]{%
    \def\@lesson{Lecture #1 (#2)}%
    \subsection*{\@lesson}%
}

\newcommand\mymaketitle{%
  \begin{titlepage}
    \null\vfil\vskip 40\p@
    \begin{center}
      {\LARGE \@title \par}
      \vskip 2.5em
      {\large \lineskip .75em \@author \par}
      \vskip 1.5em
      {\large \@date \par}
    \end{center}\par
    \@thanks
    \vfil\null
  \end{titlepage}
}


\definecolor{codegreen}{rgb}{0,0.6,0}
\definecolor{codegray}{rgb}{0.5,0.5,0.5}
\definecolor{codepurple}{rgb}{0.58,0,0.82}

\lstset{
    language=Python,
    basicstyle=\ttfamily\small,
    keywordstyle=\color{blue},
    commentstyle=\color{codegreen},
    stringstyle=\color{codepurple},
    breaklines=true,
    numbers=left,
    numberstyle=\tiny\color{codegray},
    showstringspaces=false
}
